% !TEX TS-program = pdflatex
% !TEX encoding = UTF-8 Unicode

% This is a simple template for a LaTeX document using the "article" class.
% See "book", "report", "letter" for other types of document.

\documentclass[12pt]{report} % use larger type; default would be 10pt

\usepackage[utf8]{inputenc} % set input encoding (not needed with XeLaTeX)

%%% Examples of Article customizations
% These packages are optional, depending whether you want the features they provide.
% See the LaTeX Companion or other references for full information.

%%% PAGE DIMENSIONS
\usepackage{geometry} % to change the page dimensions
\geometry{a4paper} % or letterpaper (US) or a5paper or....
% \geometry{margin=2in} % for example, change the margins to 2 inches all round
% \geometry{landscape} % set up the page for landscape
%   read geometry.pdf for detailed page layout information

\usepackage{graphicx} % support the \includegraphics command and options

% \usepackage[parfill]{parskip} % Activate to begin paragraphs with an empty line rather than an indent

%%% PACKAGES
\usepackage{booktabs} % for much better looking tables
\usepackage{array} % for better arrays (eg matrices) in maths
\usepackage{paralist} % very flexible & customisable lists (eg. enumerate/itemize, etc.)
\usepackage{verbatim} % adds environment for commenting out blocks of text & for better verbatim
\usepackage[font=small,labelfont=it]{caption}
\usepackage{subfig} % make it possible to include more than one captioned figure/table in a single float
\usepackage{pdfpages}
\usepackage{hyperref}
\usepackage{bibtopic}
\usepackage{wrapfig}
% These packages are all incorporated in the memoir class to one degree or another...

%%% HEADERS & FOOTERS
\usepackage{fancyhdr} % This should be set AFTER setting up the page geometry
\pagestyle{fancy} % options: empty , plain , fancy
\renewcommand{\headrulewidth}{0pt} % customise the layout...
\lhead{}\chead{}\rhead{}
\lfoot{}\cfoot{\thepage}\rfoot{}

%%% SECTION TITLE APPEARANCE
\usepackage{sectsty}
\allsectionsfont{\sffamily\mdseries\upshape} % (See the fntguide.pdf for font help)
\setcounter{secnumdepth}{-2} 
% (This matches ConTeXt defaults)

%%% RULE

\newcommand{\HRule}{\rule{\linewidth}{0.5mm}}

%%% BIBLIOGRAPHY

\usepackage[sectionbib]{apacite}                           %bibliography in apa-style

%%% ToC (table of contents) APPEARANCE
\usepackage[titles,subfigure]{tocloft} % Alter the style of the Table of Contents
\renewcommand{\cftsecfont}{\rmfamily\mdseries\upshape}
\renewcommand{\cftsecpagefont}{\rmfamily\mdseries\upshape} % No bold!

%%% END Article customizations

%%% The "real" document content comes below...

\begin{document}

\begin{titlepage}

\begin{center}


% Upper part of the page
\includegraphics[width=1\textwidth]{./logo}\\[1cm]    

\textsc{\Large Bachelor Thesis}\\[0.5cm]
\textsc{\Large {[}201000166{]}}\\[0.5cm]


% Title
\HRule \\[0.4cm]
{ \huge \bfseries Research Proposal}\\[0.4cm]

\HRule \\[1.5cm]

% Author and supervisor
\begin{minipage}{0.4\textwidth}
\begin{flushleft} \large
\emph{Author:}\\
Micha \textsc{van den Enk} \\
{[}s1004654{]} \\
\end{flushleft}
\end{minipage}
\begin{minipage}{0.4\textwidth}
\begin{flushright} \large
\emph{Supervisors:} \\
H. H. \textsc{Leemkuil} \\
\end{flushright}
\end{minipage}

\vfill

% Bottom of the page
{\large DATE}

\end{center}

\end{titlepage}

\tableofcontents

\bibliographystyle{apacite}

\chapter{Preface}

In this document the reader can find a proposal for designing a course on quantum mechanics in a qCraft learning environment. This is an assignment executed for a bachelor thesis. The document contains a table with general information, a short summary of the assignment, a detailed description of the assignment wth the rationale, the conceptual framework and the relevance, the design of the formative evaluation and a planning.

\chapter{General Information}

\begin{center}
\begin{tabular}{ l | p{8cm} }
Researcher & Micha van den Enk (s1004654) \\ \hline
Study & Onderwijskunde \\ \hline
Study Department & Instructional Technology \\ \hline
Date & \\ \hline
First supervisor & Dr H. H. Leemkuil \\ \hline
Second supervisor & \\ \hline
Keywords & Quantum mechanics, Middle school Education, Netherlands \\ \hline
Title & \\
\end{tabular}
\end{center}

\chapter{Summary}



\chapter{Description}

\section{Rationale}

\section{Conceptual Framework}

\section{Relevance}

\chapter{Design approach}

\chapter{Planning}

\begin{center}
\begin{tabular}{ l | l }
Analyses & 24 April \\ \hline
Literature research & 8 May \\ \hline
Design & 15 May \\ \hline
Development & 22 May \\ \hline
Evaluation & 5 June \\ \hline
Conclusion/Discussion & 12 June \\ \hline
Presentation & 19 June \\ \hline
\end{tabular}
\end{center}

\begin{btSect}{references}
\btPrintCited
\end{btSect}

\end{document}
